\section{Introduction}

Damage variability in Final Fantasy XIV is often due to hit type randomness, where critical, direct, or critical-direct hit results in increased damage dealt. A random $\pm 5\%$ damage rolls is also applied. These forms of damage variability are well-understood and there exist many tools to quantify it at varying degrees of accuracy. A few jobs, Bard and Dancer, experience additional damage variability from procs, which randomly grant extra uses of damaging abilities or enable higher-potency weaponskills. While both the statistical principles governing procs and methodologies for computing damage distribution for variable action usages are generally understood, their specific applications remain underexplored.

This paper aims to (i) develop methods to compute damage distributions that account for proc variability for Bard and (ii) compare the magnitude of variability procs introduce into damage distributions compared to hit type variability alone. Since damage variance is additive, it is useful to know if this is another type of variability that should be accounted for,\footnote{It is emphasized that improved rotations, more consistent self-play, and synergistic team play will vastly outweigh all types of variability in most cases.} or if it is mainly a mechanic to make Physical Ranged rotations more engaging. Knowing its relative magnitude is useful from both a speed kill and parsing perspective for knowing how likely faster/higher damage runs are. A better understanding can also help inform if this is something worthwhile for damage calculation tools to account for.

This paper also is not meant to model Bard's rotation with pitch-perfect precision. Example calculations are not performed on rotations that meticulously account for small, but important, details like application delays and animation locks. Rather, it provides a general methodology for quantifying each source of proc variability, which can later be extended to more precise rotation models. This analysis assumes a high-level of gameplay, \textit{e.g.}, no unintentional clipping, minimal ability drift and party buff drift, and no proc over-capping. When necessary, judiciously chosen approximations will be made to help cull the combinatoric complexity that often arises with these sorts of problems.
